\subsection{Трансформація дизайн-освіти в еру ШІ: від виробника до куратора}\label{subsec:1.2}

Аналіз педагогічної категорії інформаційно-комунікативної компетентності (ІКК), здійснений у підрозділі~\ref{subsec:1.1}, засвідчив, що формування ІКК майбутніх дизайнерів потребує врахування сучасних технологічних трансформацій професійної діяльності. Серед таких трансформацій найбільш радикальною є стрімке поширення генеративного штучного інтелекту (ГШІ), що суттєво змінює як зміст дизайнерської праці, так і вимоги до професійної підготовки фахівців. Для обґрунтування структури та змісту ІКК необхідно з'ясувати, яким чином ГШІ впливає на дизайн-освіту у глобальному масштабі, які педагогічні підходи, результати та виклики фіксуються в сучасній науковій літературі. З цією метою нами було проведено скопінг-огляд (scoping review) наукових публікацій щодо застосування ГШІ в дизайн-освіті~\cite{Musiienko2026ScopingReview}, результати якого покладено в основу цього підрозділу.

Скопінг-огляд здійснено відповідно до методологічної рамки \citet{Arksey2005}, удосконаленої \citet{Levac2010} та Joanna Briggs Institute~\cite{Peters2020}, з дотриманням рекомендацій PRISMA-ScR~\cite{Tricco2018}. Пошук проведено у базі Web of Science Core Collection за 18 пошуковими запитами, що поєднували терміни дизайн-освіти (architecture, fashion design, graphic design, product design, UX design, interaction design, design education, design pedagogy, design studio) та штучного інтелекту (artificial intelligence, generative AI, ChatGPT, machine learning, deep learning, large language model). Хронологічний діапазон пошуку --- 2017--2026~рр. З 1800 ідентифікованих записів після дворівневого відбору (скринінг за назвою та анотацією, повнотекстова оцінка відповідності) до остаточної вибірки увійшло 156 досліджень~\cite{Musiienko2026ScopingReview}.

Результати скопінг-огляду систематизовано за чотирма дослідницькими запитаннями: (1) масштаб, характер і природа досліджень ШІ у дизайн-освіті; (2) педагогічні підходи, рамки та стратегії впровадження; (3) результати, переваги, виклики та етичні аспекти; (4) лакуни в поточній літературі та пріоритетні напрями подальших досліджень. Нижче наведено детальний аналіз отриманих даних, переосмислених крізь призму формування ІКК майбутніх дизайнерів.

\subsubsection{Темпоральний розподіл публікацій та динаміка дослідницького інтересу}\label{sssec:1.2.1}

Аналіз хронологічного розподілу 156 включених досліджень виявив безпрецедентне зростання дослідницької активності, безпосередньо пов'язане з появою ChatGPT у листопаді 2022~р. До цієї дати (період 2017--2022~рр.) було опубліковано лише 21 дослідження (13,5\,\%), тоді як після неї --- 135 досліджень (86,5\,\%), що становить 6,4-кратне зростання~\cite{Musiienko2026ScopingReview}. Пік публікаційної активності припадає на 2025~р. --- 87 досліджень (55,8\,\% від загальної кількості), що свідчить про стійке прискорення дослідницького інтересу. Візуалізацію темпорального розподілу наведено на рис.~\ref{fig:1.2_temporal}.

\begin{figure}[htbp]
\centering
\begin{tikzpicture}
\begin{axis}[
    ybar,
    width=0.88\textwidth,
    height=7cm,
    bar width=0.55cm,
    ylabel={Кількість досліджень},
    xlabel={Рік публікації},
    symbolic x coords={2018, 2019, 2020, 2021, 2022, 2023, 2024, 2025, 2026},
    xtick=data,
    ymin=0,
    ymax=100,
    ytick={0, 20, 40, 60, 80, 100},
    nodes near coords,
    nodes near coords align={vertical},
    every node near coord/.append style={
        font=\footnotesize,
        /pgf/number format/assume math mode=true,
    },
    visualization depends on=\thisrow{show} \as \showmark,
    every node near coord/.append style={
        opacity=\showmark
    },
    legend style={at={(0.02,0.98)}, anchor=north west, font=\small},
    area legend,
    enlarge x limits=0.1,
]

% До ChatGPT
\addplot[fill=temporalblue, draw=blue] table[meta=show] {
    x       y   show
    2018    2   1
    2019    2   1
    2020    4   1
    2021    9   1
    2022    4   1
    2023    0   0
    2024    0   0
    2025    0   0
    2026    0   0
};

% Після ChatGPT
\addplot[fill=temporalred, draw=red] table[meta=show] {
    x       y   show
    2018    0   0
    2019    0   0
    2020    0   0
    2021    0   0
    2022    0   0
    2023    9   1
    2024    36  1
    2025    87  1
    2026    3   1
};

\legend{До ChatGPT (\textit{n}=21), Після ChatGPT (\textit{n}=135)}

% Вертикальна лінія між 2022 і 2023
\draw[dashed, thick, gray] (rel axis cs:0.555,0) -- (rel axis cs:0.555,0.95);
\node[font=\footnotesize, align=center] at (rel axis cs:0.555,0.92) {Реліз\\ChatGPT};

\end{axis}
\end{tikzpicture}
\caption{Темпоральний розподіл досліджень ШІ у дизайн-освіті (6,4-кратне зростання після ChatGPT)}
\label{fig:1.2_temporal}
\end{figure}

Такий вибуховий характер зростання публікаційної активності є показовим не лише з бібліометричної точки зору, а й у контексті формування ІКК. Стрімкість технологічних змін означає, що освітні програми підготовки дизайнерів опинилися в ситуації, коли швидкість впровадження ГШІ суттєво випереджає розвиток відповідних педагогічних стратегій та навчальних методик. Як зазначає \citet{Fleischmann2024}, генеративний ШІ вже стає незамінним інструментом для ідеації та прототипування --- двох фундаментальних навичок, що формуються у студійній педагогіці дизайну. Водночас дані опитувань свідчать, що майже 85\,\% викладачів і 86\,\% студентів використовували ШІ протягом 2024--2025 навчального року, проте відсутність системного навчання створила так званий <<компетентнісний розрив>>, за якого швидкість адаптації технологій випереджає розвиток критичної грамотності та етичних рамок~\cite{Pudasaini2025}.

Ця ситуація безпосередньо стосується проблеми формування ІКК: інформаційно-аналітичний компонент ІКК має включати здатність критично оцінювати можливості та обмеження ГШІ, технологічний компонент --- практичні вміння ефективного використання ШІ-інструментів, а рефлексивний компонент --- усвідомлення етичних та правових аспектів їхнього застосування у професійній дизайнерській діяльності.

\subsubsection{Географічний розподіл досліджень та глобальний контекст}\label{sssec:1.2.2}

Географічний аналіз за афіліацією першого автора виявив значну нерівномірність дослідницької активності. Домінуючу позицію займає Китай --- 49 досліджень (31,4\,\%), за ним слідують США --- 21 дослідження (13,5\,\%), Південна Корея --- 9 (5,8\,\%), Тайвань --- 8 (5,1\,\%), Велика Британія --- 7 (4,5\,\%), Німеччина та Туреччина --- по 6 (3,8\,\%), Австралія --- 5 (3,2\,\%), Іспанія --- 4 (2,6\,\%), Італія --- 3 (1,9\,\%), решта країн --- 38 (24,4\,\%)~\cite{Musiienko2026ScopingReview}.

Регіональний аналіз підтверджує концентрацію досліджень у Східній Азії: Китай, Тайвань і Південна Корея разом становлять 42,3\,\% усіх публікацій. Європа (включаючи Велику Британію, Німеччину, Туреччину, Іспанію, Італію) забезпечує 16,7\,\%, Північна Америка (США) --- 13,5\,\%. Показовими є суттєві географічні лакуни: Африка, Латинська Америка та більша частина Південної Азії практично не представлені у дослідженнях, що ставить під сумнів глобальну генералізованість наявних висновків.

Для українського контексту формування ІКК цей факт має подвійне значення. По-перше, відсутність досліджень із пострадянського та центральноєвропейського простору означає, що педагогічні моделі інтеграції ШІ в дизайн-освіту потребують адаптації до специфіки національної освітньої системи. По-друге, домінування англомовних досліджень з Китаю та США вказує на необхідність ретельного культурного контекстуалізування при перенесенні зарубіжного досвіду до вітчизняної практики підготовки дизайнерів.

\subsubsection{Дисциплінарний розподіл та проблема графічного дизайну}\label{sssec:1.2.3}

Розподіл за дизайнерськими дисциплінами засвідчив домінування загальної дизайн-освіти --- 57,7\,\% досліджень, що відображає як міждисциплінарний характер дискусій щодо ШІ, так і потенційну нечіткість у дисциплінарній ідентифікації досліджень. Архітектура та промисловий/продуктовий дизайн представлені відносно рівномірно (по 12,2\,\%), що пояснюється тривалою традицією залучення комп'ютерних інструментів у ці галузі. Дизайн одягу становить 8,3\,\%, UX/UI та дизайн взаємодії --- 5,1\,\%, дизайн інтер'єру --- 1,9\,\%~\cite{Musiienko2026ScopingReview}.

Принципово важливим є критично низький рівень представленості графічного/візуального дизайну --- лише 1,9\,\% досліджень, що є парадоксальним з огляду на те, що саме інструменти генерації зображень (DALL-E, Midjourney, Stable Diffusion) є найбільш поширеними та обговорюваними серед ГШІ-технологій. Ця лакуна є особливо значущою для нашого дослідження, оскільки спеціальність 022 <<Дизайн>> в Україні охоплює графічний дизайн як одну з основних спеціалізацій. Відсутність достатньої дослідницької бази щодо впливу ГШІ на підготовку графічних дизайнерів підкреслює актуальність розробки вітчизняної моделі формування ІКК, яка б враховувала специфіку саме цієї дизайнерської дисципліни.

\subsubsection{Методологічний ландшафт та якість досліджень}\label{sssec:1.2.4}

Аналіз методологічних підходів включених досліджень виявив суттєві проблеми з науковою строгістю наявної літератури. Концептуальні та теоретичні роботи становлять найбільшу категорію --- 36 досліджень (23,1\,\%), за ними йдуть експериментальні дослідження --- 35 (22,4\,\%), емпіричні без чіткої методології --- 31 (19,9\,\%), огляди --- 17 (10,9\,\%), якісні дослідження та опитування --- по 14 (9,0\,\%), кейс-стаді --- 7 (4,5\,\%), дизайн-дослідження --- 1 (0,6\,\%), змішані методи --- 1 (0,6\,\%)~\cite{Musiienko2026ScopingReview}.

За рівнем якості доказів: дослідження вищої якості (експериментальні, дизайн-дослідження, змішані методи) становлять лише 23,7\,\%; середньої якості (опитування, якісні, кейс-стаді) --- 22,4\,\%; нижчої якості (концептуальні, емпіричні без методології, огляди) --- 53,8\,\%. Такий розподіл свідчить про обмежену емпіричну обґрунтованість поточної літератури. Відсутність лонгітюдних досліджень, що відстежували б довготривалий вплив ГШІ на навчальні результати, є особливо проблематичною для обґрунтування освітніх реформ.

У контексті формування ІКК методологічна незрілість поля досліджень означає, що педагогічні рішення щодо інтеграції ШІ у підготовку дизайнерів базуються переважно на теоретичних міркуваннях та фрагментарному емпіричному досвіді, а не на системній доказовій базі. Це додатково обґрунтовує необхідність розробки науково виваженої моделі формування ІКК, заснованої на комплексному аналізі як міжнародного, так і вітчизняного досвіду.

\subsubsection{Технологічний фокус досліджень}\label{sssec:1.2.5}

Аналіз технологічного фокусу засвідчив, що значна частина досліджень оперує загальними термінами без конкретизації конкретних технологій. ШІ загалом (AI general) є предметом 55 досліджень (35,3\,\%), генеративний ШІ загалом --- 29 (18,6\,\%), великі мовні моделі (LLM)/ChatGPT --- 27 (17,3\,\%), машинне навчання/глибоке навчання --- 24 (15,4\,\%), великі мовні моделі (не конкретизовані) --- 12 (7,7\,\%), ШІ генерації зображень --- 9 (5,8\,\%)~\cite{Musiienko2026ScopingReview}.

Лише 36 досліджень (23,1\,\%) згадували конкретні ШІ-інструменти за назвою. Серед них: ChatGPT --- 23 дослідження (14,7\,\%), Midjourney --- 7 (4,5\,\%), Copilot --- 6 (3,8\,\%), GPT-4 --- 4 (2,6\,\%). \citet{Turchi2023} надали один із небагатьох детальних аналізів впливу конкретних інструментів (DALL-E, Midjourney, Stable Diffusion) на творчі процеси, виявивши, що <<взаємодія з інструментами генерації зображень може впливати на креативність>> через зворотні зв'язки, що формують ідеаційні патерни користувача.

Для технологічного компоненту ІКК цей аналіз вказує на необхідність формування у студентів-дизайнерів не прив'язаних до конкретних інструментів, а загальних (трансферабельних) технологічних умінь: здатності оцінювати та обирати відповідні ШІ-інструменти для конкретних проєктних завдань, адаптуватися до швидкої зміни технологічного ландшафту, критично аналізувати можливості та обмеження різних платформ.

\subsubsection{Тематичні кластери: педагогічні фокуси досліджень}\label{sssec:1.2.6}

Тематичний аналіз педагогічних фокусів виявив дев'ять основних напрямів, що відображають ключові проблеми інтеграції ШІ в дизайн-освіту (дослідження могли мати кілька тематичних фокусів, тому відсотки сумуються до більш ніж 100\,\%). Візуалізацію тематичних кластерів наведено на рис.~\ref{fig:1.2_themes}.

\begin{figure}[htbp]
\centering
\begin{tikzpicture}
\newcommand{\barscale}{0.23}
\newcommand{\barht}{0.35}
\newcommand{\ygap}{0.65}
% Labels
\foreach \lbl/\val/\idx in {%
  Розвиток навичок/7.1/0,%
  Критичне мислення/9.0/1,%
  Промпт-інженерія/9.0/2,%
  Колаборація/10.3/3,%
  Етика та доброчесність/12.8/4,%
  Студійне навчання/12.8/5,%
  Дизайн навч.\ програм/22.4/6,%
  Оцінювання/27.6/7,%
  Розвиток креативності/35.9/8%
}{%
  \pgfmathsetmacro{\ypos}{\idx*\ygap}%
  \draw[fill=themebarfill, draw=themebaredge] (0,\ypos) rectangle (\val*\barscale, \ypos+\barht);
  \node[anchor=east, font=\small] at (-0.15,\ypos+\barht/2) {\lbl};
  \node[anchor=west, font=\footnotesize] at (\val*\barscale+0.1, \ypos+\barht/2) {\val\%};
}
% X axis
\draw[->] (0,-0.2) -- (10,{-0.2});
\foreach \x/\lab in {0/0, 2.3/10, 4.6/20, 6.9/30, 9.2/40} {
  \draw (\x,-0.3) -- (\x,-0.1);
  \node[font=\footnotesize] at (\x,-0.55) {\lab};
}
\node[font=\small] at (4.6,-1) {\% досліджень};
\end{tikzpicture}
\caption{Розподіл педагогічних тематичних фокусів у дослідженнях ШІ в дизайн-освіті}
\label{fig:1.2_themes}
\end{figure}

\textbf{Розвиток креативності} є домінуючим педагогічним фокусом (35,9\,\%), що відображає фундаментальну зацікавленість дизайн-освіти у збереженні та розвитку творчого потенціалу в умовах технологічних трансформацій. \citet{MedelVera2025} встановили, що ГШІ може функціонувати <<як колаборатор і провокатор у дизайн-педагогіці, сприяючи творчій ідеації та водночас спонукаючи до нових педагогічних стратегій, зосереджених на промпт-грамотності та рефлексивному дизайні>>. Семантична насиченість текстових промптів корелює з більш відкритими та виразними ШІ-генерованими зображеннями, що свідчить про те, що промпт-грамотність є справжньою новою формою дизайнерської компетентності, а не лише технічним навиком.

У контексті ІКК ці знахідки обґрунтовують необхідність включення до інформаційно-аналітичного компоненту здатності формулювати ефективні промпти (текстові запити до ШІ-систем), критично оцінювати згенеровані результати та ітеративно вдосконалювати запити для досягнення бажаного творчого результату.

\textbf{Оцінювання} є другим за поширеністю тематичним фокусом (27,6\,\%), що відображає глибоку стурбованість академічної спільноти питаннями валідності оцінювання навчальних досягнень в умовах, коли ШІ-допомога є повсюдно доступною. \citet{Perkins2023} продемонстрував, що наявні інструменти виявлення плагіату <<ще не оснащені методами точного виявлення ШІ-генерованого тексту>>, що принципово ускладнює традиційну систему оцінювання, засновану на кінцевому артефакті як показнику індивідуального навчання.

\citet{Deluca2025} запропонували модель AI3, що інтегрує ШІ з академічною доброчесністю та інноваціями в оцінюванні. Ця модель передбачає доповнення традиційних методів оцінювання фінальних артефактів процесною документацією, що розкриває внесок людини; усними опитуваннями, що перевіряють розуміння та рефлексію прийнятих рішень; аудиторними вправами для встановлення базового рівня компетентностей; портфоліо-кураторством з акцентом на відбір і рефлексію; та оцінюванням навичок критичного аналізу дизайну.

Для рефлексивного компоненту ІКК ці знахідки є принципово важливими: формування ІКК має включати розвиток здатності студентів-дизайнерів до критичної самооцінки власного творчого внеску у проєкт, розмежування результатів власної інтелектуальної роботи та ШІ-генерованого контенту, а також рефлексії щодо ролі ШІ-інструментів у власному творчому процесі.

\textbf{Дизайн навчальних програм} є третім за значущістю фокусом (22,4\,\%), що свідчить про усвідомлення необхідності системних змін у змісті та структурі дизайн-освіти. Цей напрям безпосередньо пов'язаний із завданнями нашого дослідження щодо обґрунтування організаційно-педагогічних умов формування ІКК.

\textbf{Студійне та проєктне навчання} (12,8\,\%) та \textbf{етика і доброчесність} (12,8\,\%) посідають рівноцінні позиції, що відображає одночасну увагу до практичних аспектів інтеграції ШІ у студійну педагогіку та до нормативних рамок такої інтеграції. \citet{Iranmanesh2024} наводять аргументи як на користь інтеграції ШІ в архітектурну освіту --- для підготовки студентів до сучасної професійної практики, так і проти --- через загрозу підриву фундаментальної педагогіки <<навчання через діяльність>> (learning by doing).

\textbf{Колаборація} (10,3\,\%), \textbf{промпт-інженерія} та \textbf{критичне мислення} (по 9,0\,\%) виокремлюються як нові компетентнісні зони. Промпт-інженерія є безпосередньо новим компонентом технологічної складової ІКК. \citet{Lee2024} розробили <<ШІ-промпт-путівник>> для освіти з дизайну одягу, виявивши, що студенти оцінили його у 4,7 з 5 балів за корисність для генерування творчих дизайнерських результатів через платформи ChatGPT та Midjourney.

\textbf{Розвиток навичок} (7,1\,\%) є найменш представленим фокусом, що є дещо парадоксальним з огляду на фундаментальні зміни у переліку необхідних дизайнерських умінь. Саме ця лакуна обґрунтовує необхідність детального дослідження структури та змісту навичок, що мають формуватися у рамках ІКК майбутніх дизайнерів.

\subsubsection{Теоретична недорозвиненість поля досліджень}\label{sssec:1.2.7}

Одним із найбільш значущих результатів скопінг-огляду є виявлення суттєвої теоретичної лакуни: лише 41 дослідження (26,3\,\%) використовувало явні теоретичні рамки, тоді як 115 досліджень (73,7\,\%) є атеоретичними~\cite{Musiienko2026ScopingReview}. Розподіл теоретичних рамок серед досліджень, що їх використовували, наведено у табл.~\ref{tab:1.2_theory}.

\begin{table}[htbp]
\centering
\caption{Теоретичні рамки, використані у дослідженнях ШІ в дизайн-освіті}
\label{tab:1.2_theory}
\small
\begin{tabular}{@{}lccp{6cm}@{}}
\toprule
\textbf{Теоретична рамка} & \textbf{\textit{n}} & \textbf{\%} & \textbf{Типове застосування} \\
\midrule
Не визначена & 115 & 73,7 & --- \\
Дизайн-мислення & 11 & 7,1 & Структурування фаз ШІ-асистованої ідеації \\
Проєктне навчання & 8 & 5,1 & Інтеграція ШІ у студійні проєкти \\
Досвідне навчання & 7 & 4,5 & Практичне дослідження ШІ-інструментів \\
TPACK & 6 & 3,8 & Інтеграція технологій, педагогіки та контенту \\
UTAUT & 4 & 2,6 & Фактори прийняття технологій \\
Теорія когнітивного навантаження & 3 & 1,9 & Вплив ШІ на ментальне зусилля \\
Конструктивізм & 3 & 1,9 & Студентоцентроване навчання з ШІ \\
\bottomrule
\end{tabular}
\end{table}

Домінування атеоретичних досліджень (73,7\,\%) є найбільш проблемним висновком огляду. Без теоретичного обґрунтування дослідження продукують спостереження без пояснень, накопичують факти без побудови системного розуміння. Механізми, через які ШІ впливає на навчання, креативність або дизайнерську когніцію, залишаються недостатньо специфікованими.

Серед використаних теоретичних рамок особливо релевантною для формування ІКК є рамка TPACK (Technological Pedagogical Content Knowledge --- технологічно-педагогічне знання контенту). \citet{Lee2024} застосували TPACK до освіти з дизайну одягу, встановивши, що інтеграція ШІ <<підвищує креативність та ефективність студентів>>, коли технологічне, педагогічне та контентне знання систематично узгоджені. \citet{Bower2017} обґрунтовує, що рамка TPACK є продуктивним інструментом для розуміння того, як викладачі можуть ефективно інтегрувати ШІ-інструменти у навчальний процес. \citet{Boschman2015} демонструють, що TPACK-компетентність викладачів розвивається через спільну проєктну діяльність, що вказує на шляхи професійного розвитку для інтеграції ШІ.

\citet{Polly2020} аргументують, що зона найближчого розвитку Л.~С.~Виготського може інформувати розвиток TPACK, пропонуючи поступовий шлях формування компетентностей інтеграції ШІ. \citet{Dong2025} застосували теорію когнітивного навантаження та структурне моделювання, встановивши, що <<ШІ-асистовані шляхи суттєво підвищують якість дизайну (72,43 проти 65,60 у традиційних шляхах)>> шляхом зниження побічного когнітивного навантаження при одночасному підвищенні інвестицій у продуктивне когнітивне навантаження. Це свідчить про те, що зважене використання ШІ може фактично посилювати, а не підривати глибоке інноваційне мислення студентів.

\citet{Liu2025} використали як UTAUT (Unified Theory of Acceptance and Use of Technology --- уніфіковану теорію прийняття та використання технологій), так і AHP (Analytic Hierarchy Process --- метод аналізу ієрархій) для оцінювання ГШІ-інструментів у дизайні інтер'єру, встановивши, що <<Stable Diffusion перевершує у креативності, тоді як Midjourney домінує в естетиці та функціональності>>. Цей порівняльний підхід представляє новий методологічний напрям для обґрунтованого вибору інструментів.

Для обґрунтування моделі формування ІКК теоретична лакуна одночасно є і проблемою, і можливістю: вона засвідчує необхідність розробки теоретично обґрунтованого підходу, що інтегрував би педагогічні, технологічні та дизайнерські виміри формування компетентності. Рамка TPACK, доповнена рефлексивним компонентом і специфікою дизайнерської діяльності, може стати одним із таких теоретичних фундаментів.

\subsubsection{Результати впровадження ШІ: переваги та обмеження}\label{sssec:1.2.8}

Суттєвою знахідкою огляду є те, що 66,7\,\% досліджень не зафіксували конкретних вимірюваних результатів, що суттєво обмежує доказову базу ефективності ШІ-інтеграції~\cite{Musiienko2026ScopingReview}. Серед досліджень, що звітували про результати: навчальні досягнення --- 30 (19,2\,\%), прийняття технологій --- 20 (12,8\,\%), підвищення залученості --- 6 (3,8\,\%), покращення креативності --- 5 (3,2\,\%), зростання ефективності --- 1 (0,6\,\%).

Серед досліджень із зафіксованими результатами виділяються роботи, що використовували більш строгі методологічні підходи. \citet{MedelVera2025} застосували Creativity Support Index (CSI) у поєднанні із семантичним аналізом та естетичним оцінюванням, встановивши <<сильну підтримку дослідницької та цілеспрямованої креативності, з високими показниками за вимірами дослідження та цінності результатів>>. \citet{Dong2025} за допомогою когнітивної теорії навантаження та структурного моделювання виявили, що ШІ-асистовані шляхи знижують побічне когнітивне навантаження з 0,907 до 0,017, водночас підвищуючи інвестиції у продуктивне когнітивне навантаження, що сприяє глибокому інноваційному мисленню.

Аналогічно, 74,4\,\% досліджень не задокументували конкретних викликів. Серед зафіксованих: упередженість/справедливість (bias/fairness) --- 15 (9,6\,\%), впроваджувальні виклики --- 10 (6,4\,\%), потреби у підвищенні кваліфікації викладачів --- 6 (3,8\,\%), академічна доброчесність --- 6 (3,8\,\%), авторське право та інтелектуальна власність --- 3 (1,9\,\%), надмірна залежність --- 2 (1,3\,\%)~\cite{Musiienko2026ScopingReview}.

Таке суттєве недозвітування як результатів (66,7\,\%), так і викликів (74,4\,\%) є ключовим обмеженням поточного стану досліджень. Для формування ІКК це означає, що ефективність конкретних педагогічних стратегій інтеграції ШІ в дизайн-освіту залишається переважно бездоказовою, що підкреслює необхідність проведення власного емпіричного дослідження у рамках даної дисертаційної роботи.

\subsubsection{Упередженість ШІ, етика та академічна доброчесність}\label{sssec:1.2.9}

Серед задокументованих викликів упередженість ШІ (9,6\,\%) є найбільш дослідженою проблемою. \citet{Currie2024} провели систематичну оцінку упередженості генеративних ШІ text-to-image інструментів, виявивши, що серед згенерованих зображень <<60,6\,\% зображали чоловіків, а 87,7\,\% мали світлий тон шкіри>>. При цьому DALL-E 3 містить <<вбудовані упередження, пов'язані з гендером та етнічною приналежністю, що вимагають більш критичного оцінювання>>. Ці дані мають безпосередні наслідки для дизайн-освіти, де репрезентативність і культурна чутливість є важливими професійними цінностями.

Для інформаційно-аналітичного компоненту ІКК ці знахідки обґрунтовують необхідність формування у майбутніх дизайнерів здатності розпізнавати та протидіяти алгоритмічній упередженості у ШІ-генерованому контенті. Дизайнер, який використовує ШІ-інструменти без критичного аналізу їхніх результатів, ризикує відтворювати та посилювати існуючі стереотипи у своїй професійній роботі.

Проблеми академічної доброчесності зафіксовано лише у 3,8\,\% досліджень, що є дивно низьким показником з огляду на інтенсивність публічного дискурсу з цієї теми. \citet{Perkins2023} надає комплексний аналіз, аргументуючи, що <<визначальним є не саме використання ШІ-інструментів студентом, а те, чи зроблено таке використання прозорим>>. \citet{Pudasaini2025} у своєму огляді методів виявлення ШІ-генерованого плагіату доходять висновку, що наявні інструменти <<не здатні надійно розрізняти людські та ШІ-написані тексти>>.

Питання авторського права та інтелектуальної власності досліджено у 1,9\,\% робіт. \citet{Mazzi2024} аналізує проблему авторства у ШІ-генерованих творах, досліджуючи оригінальність текстових промптів та їхню кореляцію з результатами. Ці питання набувають особливої гостроти для дизайнерської професії, де оригінальність і авторство є ключовими категоріями.

\citet{Buragohain2025} досліджують етичні та педагогічні перспективи інтеграції ChatGPT у вищу освіту країн ASEAN, ідентифікуючи ключові етичні проблеми, пов'язані з конфіденційністю даних, алгоритмічною упередженістю та впливом на навчальні практики. Рекомендації авторів щодо відповідального використання ГШІ є релевантними для формування етичного виміру ІКК.

Важливим нормативним контекстом для етичного виміру ІКК є Закон України <<Про штучний інтелект>> (№\,4742-IX від 15~лютого 2025~р.)~\cite{ZakonAI2025}. Стаття~8(4) Закону встановлює, що ШІ-генерований контент, використаний в академічних роботах, має бути розкритий із зазначенням методології його отримання. Стаття~29 кваліфікує подання ШІ-генерованих результатів як власних без належного розкриття як порушення академічної доброчесності. Ці положення формують законодавчу рамку для рефлексивного компоненту ІКК: майбутній дизайнер має не лише вміти критично оцінювати ШІ-генерований контент, а й дотримуватися правових вимог щодо прозорого звітування про використання ШІ у професійній та академічній діяльності.

У сукупності аналіз етичних аспектів обґрунтовує включення до рефлексивного компоненту ІКК здатності до етичної рефлексії щодо: (а) упередженості ШІ-генерованого контенту та стратегій її мінімізації; (б) академічної доброчесності --- прозорого звітування про роль ШІ у творчому процесі; (в) авторського права --- відповідального використання ШІ-інструментів з дотриманням правових норм; (г)~дотримання вимог Закону України <<Про штучний інтелект>> щодо розкриття ШІ-генерованого контенту~\cite{ZakonAI2025}.

\subsubsection{Парадигмальний зсув: від виробника до куратора}\label{sssec:1.2.10}

Наскрізною темою проаналізованої літератури є трансформація ролі дизайнера від виробника (maker) до куратора (curator) або оркестратора ШІ-генерованих результатів~\cite{Musiienko2026ScopingReview}. Ця трансформація паралельна попереднім технологічним переходам --- від ремісника до оператора САПР, від креслярника до параметричного дизайнера --- але з принциповою відмінністю: генеративний ШІ втручається не лише у виконання, а й безпосередньо в ідеацію.

\citet{Schon1983} обґрунтував, що <<рефлексія у дії>> (reflection-in-action) --- інтеграція мислення та діяння через усвідомлену практику --- є фундаментальним механізмом розвитку дизайнерської експертності. Дизайн-освіта історично привілейовувала <<навчання через діяльність>> (learning by doing) --- ітеративний розвиток дизайнерського судження через практичну роботу з матеріалами, процесами та критикою. Студійна модель, описана \citet{Shulman2005} як одна з <<фірмових педагогік>> професій, створює соціальний та матеріальний контекст, у якому студенти розвивають не лише технічні навички, а й професійну ідентичність, критичну чутливість і те, що \citet{Cross2011} називає <<дизайнерськими способами пізнання>>.

\citet{Sennett2008} аргументував, що майстерність --- тривала практика виготовлення речей якісно --- розвиває не лише вміння, а й характер, терпіння та судження. Ключове освітнє питання полягає у тому, чи може кураторство ШІ-результатів розвивати еквівалентні якості, або ж щось суттєве втрачається, коли рука більше не керує виготовленням. Обмежена увага до цього питання у поточній літературі (розвиток навичок з'являється лише у 7,1\,\% досліджень) становить значну концептуальну лакуну.

Водночас \citet{MedelVera2025} пропонують продуктивне бачення, виявивши, що ГШІ може функціонувати <<як колаборатор і провокатор у дизайн-педагогіці>>. Їхній акцент на промпт-грамотності та рефлексивному дизайні свідчить про те, що нові форми дизайнерського мислення можуть виникати з людино-ШІ-колаборації, а не бути просто зредукованими версіями традиційної практики.

Для концепції ІКК парадигмальний зсув <<від виробника до куратора>> є ключовим: якщо роль дизайнера змінюється від безпосереднього виготовлення артефактів до кураторства, оркестрування та критичного оцінювання ШІ-генерованих результатів, то й компетентнісна модель підготовки має відображати ці зміни. Інформаційно-аналітичний компонент ІКК набуває нової функції --- критичного оцінювання та відбору ШІ-генерованого контенту; комунікативний компонент --- здатності ефективно комунікувати з ШІ-системами через промпти та інтерпретувати їхні результати для професійної аудиторії; технологічний компонент --- майстерності використання ШІ-інструментів як творчих асистентів; рефлексивний компонент --- здатності до метакогнітивної рефлексії щодо власного творчого процесу в умовах людино-ШІ-колаборації.

\subsubsection{Промпт-інженерія як нова компетентність}\label{sssec:1.2.11}

Виокремлення промпт-інженерії як педагогічного фокусу (9,0\,\% досліджень) сигналізує про визнання нових компетентностей, необхідних у ШІ-опосередкованій дизайнерській роботі. Промпт-інженерія вже викладається в таких інституціях, як Georgia Tech, Columbia University та RISD~\cite{MedelVera2025}. Навчальні програми охоплюють <<основи промптування>>, ітеративний дизайн через <<ланцюжки промптів>> (prompt chaining) та критичне оцінювання ШІ-результатів, включаючи захист від <<галюцинацій>> --- ситуацій, коли ШІ генерує переконливий, але фактично некоректний контент.

\citet{MedelVera2025} демонструють, що семантична насиченість текстових промптів корелює з більш відкритими та виразними ШІ-генерованими зображеннями. Це свідчить про те, що промпт-грамотність є справжньою новою формою дизайнерської компетентності, а не просто технічним навиком. Ефективний промпт вимагає глибокого розуміння дизайнерської мови, візуальної культури, стилістичних конвенцій --- тобто саме тих знань, що формуються у процесі професійної підготовки дизайнера.

Для технологічного компоненту ІКК промпт-інженерія є конкретною новою компетентністю, що має бути включена до змісту підготовки. Це не просто уміння формулювати текстовий запит, а комплексна здатність, що включає: (а) вербалізацію візуальних концепцій засобами природної мови; (б) ітеративне вдосконалення промптів на основі аналізу згенерованих результатів; (в) використання стратегій ланцюжкового промптування для складних дизайнерських завдань; (г) критичне оцінювання та курування ШІ-генерованого контенту; (д) інтеграцію ШІ-результатів у цілісний дизайнерський проєкт.

\subsubsection{Дизайн навчальних програм та педагогічні імплікації}\label{sssec:1.2.12}

Аналіз педагогічних фокусів дозволяє ідентифікувати п'ять ключових компетентнісних зон, що мають бути відображені у змісті дизайн-освіти~\cite{Musiienko2026ScopingReview}:

\begin{itemize}
\item \textit{ШІ-грамотність} --- критичне розуміння принципів функціонування ШІ-систем, їхніх можливостей, обмежень та суспільних наслідків~\cite{Pudasaini2025}. Ця компетентність відповідає інформаційно-аналітичному компоненту ІКК.

\item \textit{Промпт-майстерність} --- ітеративна, рефлексивна практика формулювання ефективних промптів як нова форма дизайнерського мислення~\cite{Lee2024}. Ця компетентність перетинає технологічний та комунікативний компоненти ІКК.

\item \textit{Кураторство та судження} --- здатність оцінювати, відбирати та вдосконалювати ШІ-результати відповідно до дизайнерських критеріїв~\cite{MedelVera2025}. Ця компетентність є ядром рефлексивного компоненту ІКК.

\item \textit{Людино-ШІ-колаборація} --- стратегії інтеграції ШІ у персональний творчий робочий процес зі збереженням авторського голосу~\cite{Turchi2023}. Ця компетентність поєднує комунікативний та технологічний компоненти ІКК.

\item \textit{Етичне міркування} --- рамки для прийняття рішень щодо доречності та способів використання ШІ~\cite{Buragohain2025}. Ця компетентність реалізується через рефлексивний компонент ІКК.
\end{itemize}

Зіставлення цих п'яти компетентнісних зон із чотирикомпонентною структурою ІКК (інформаційно-аналітичний, комунікативний, технологічний, рефлексивний компоненти) засвідчує їхню повну взаємовідповідність, що підтверджує обґрунтованість запропонованої структури ІКК для умов ШІ-трансформованої дизайн-освіти.

\subsubsection{Потреби у підвищенні кваліфікації викладачів}\label{sssec:1.2.13}

Потреби у професійному розвитку викладачів зафіксовано лише у 3,8\,\% досліджень, що є критично низьким показником з огляду на те, що саме викладачі є агентами педагогічних змін~\cite{Musiienko2026ScopingReview}. \citet{Fleischmann2024} виявила, що студенти використовують ГШІ <<стихійно>> зі <<скептичним ставленням до його творчих результатів>>, що свідчить про необхідність, але водночас відсутність, викладацького керівництва.

Ефективна інтеграція ШІ у дизайн-освіту вимагає від викладачів: (а) розвитку власних ШІ-компетентностей; (б) переосмислення своїх ролей від трансляторів знань до фасилітаторів навчання; (в) редизайну завдань і методів оцінювання; (г) навігації етичними складностями за відсутності чітких інституційних настанов; (д) моделювання належного використання ШІ для студентів.

Ці висновки є безпосередньо релевантними для розділу 2 нашого дослідження, де буде обґрунтовано організаційно-педагогічні умови формування ІКК, включаючи підготовку викладацького складу до роботи в умовах ШІ-трансформованої дизайн-освіти.

\subsubsection{Системний аналіз лакун у дослідженнях}\label{sssec:1.2.14}

Проведений скопінг-огляд ідентифікував шість основних лакун у поточній літературі щодо ШІ в дизайн-освіті~\cite{Musiienko2026ScopingReview}, кожна з яких має безпосередні імплікації для формування ІКК. Візуалізацію лакун наведено на рис.~\ref{fig:1.2_gaps}.

\begin{figure}[htbp]
\centering
\begin{tikzpicture}[
    gap/.style={dia/lrbox=3.2cm, minimum height=1.6cm, font=\small},
    dia/arrow/.style={-{Stealth[length=2mm]}, thick, gray}
]

% Центральний вузол
\node[ellipse, draw, thick, text width=3cm, align=center, font=\small\bfseries] (center) {Лакуни\\сучасних\\досліджень};

% Вузли лакун
\node[gap, above left=1.5cm and 1cm of center] (g1) {\textbf{Теоретична}\\73,7\,\% атеоретичних};
\node[gap, above right=1.5cm and 1cm of center] (g2) {\textbf{Методологічна}\\53,8\,\% низької якості};
\node[gap, left=2cm of center] (g3) {\textbf{Дисциплінарна}\\Графічний дизайн: 1,9\,\%};
\node[gap, right=2cm of center] (g4) {\textbf{Географічна}\\Африка/Лат. Ам.: 0\,\%};
\node[gap, below left=1.5cm and 1cm of center] (g5) {\textbf{Результативна}\\66,7\,\% без результатів};
\node[gap, below right=1.5cm and 1cm of center] (g6) {\textbf{Стейкхолдерна}\\Викладачі: 3,8\,\%};

% Стрілки
\draw[dia/arrow] (center) -- (g1);
\draw[dia/arrow] (center) -- (g2);
\draw[dia/arrow] (center) -- (g3);
\draw[dia/arrow] (center) -- (g4);
\draw[dia/arrow] (center) -- (g5);
\draw[dia/arrow] (center) -- (g6);

\end{tikzpicture}
\caption{Шість основних лакун у дослідженнях ШІ в дизайн-освіті}
\label{fig:1.2_gaps}
\end{figure}

\textit{Теоретична лакуна.} При 73,7\,\% атеоретичних досліджень поле має обмежену здатність пояснити, \textit{чому} відбуваються спостережувані явища, або будувати кумулятивне знання. Механізми, через які ШІ впливає на навчання, креативність або дизайнерську когніцію, залишаються недостатньо специфікованими. Для формування ІКК це означає необхідність розробки власної теоретичної рамки, що інтегруватиме педагогічні, технологічні та дизайнерські виміри.

\textit{Методологічна лакуна.} Переважання концептуальних робіт (23,1\,\%) та досліджень нижчої якості (53,8\,\%) обмежує обґрунтованість висновків. Відсутність лонгітюдних досліджень унеможливлює розуміння довготривалих ефектів. Лише одне дослідження використало змішані методи. Для нашого дослідження це обґрунтовує обрання комплексної методології, що поєднує кількісні та якісні підходи.

\textit{Дисциплінарна лакуна.} Графічний/візуальний дизайн (1,9\,\%) є драматично недопредставленим, незважаючи на безпосередню релевантність ШІ-генерації зображень для цієї дисципліни. Ігровий дизайн, сервіс-дизайн та текстильний/ремісничий дизайн фактично відсутні. Ця лакуна є особливо значущою для українського контексту, де графічний дизайн є однією з основних спеціалізацій у межах спеціальності 022 <<Дизайн>>.

\textit{Географічна лакуна.} Концентрація досліджень у Східній Азії (42,3\,\%) та відсутність із Африки, Латинської Америки та більшої частини Південної Азії обмежує генералізованість і виключає важливі освітні контексти. Повна відсутність досліджень з України та пострадянського простору підкреслює унікальність і актуальність даного дисертаційного дослідження.

\textit{Результативна лакуна.} При 66,7\,\% досліджень, що не специфікують вимірюваних результатів, доказова база ефективності ШІ в освіті є неповною. Дослідження, що звітують про результати, рідко використовують валідовані інструменти або порівняльні умови.

\textit{Стейкхолдерна лакуна.} Перспективи викладачів (3,8\,\%) та голоси студентів отримують обмежену увагу. Досвід, занепокоєння та стратегії основних учасників дизайн-освіти є недостатньо дослідженими.

\subsubsection{Імплікації для формування ІКК майбутніх дизайнерів}\label{sssec:1.2.15}

Узагальнення результатів скопінг-огляду 156 досліджень дозволяє сформулювати ключові висновки, що безпосередньо впливають на зміст і структуру формування ІКК майбутніх дизайнерів.

\textit{По-перше}, парадигмальний зсув <<від виробника до куратора>> вимагає переосмислення змісту всіх компонентів ІКК. Якщо традиційна дизайн-освіта акцентувала технічні навички виготовлення артефактів, то сучасна ІКК має включати навички оцінювання, відбору, рефайнування та оркестрування ШІ-генерованих результатів. Це не означає відмову від ручних умінь, а їхнє доповнення новим шаром компетентностей.

\textit{По-друге}, промпт-інженерія як нова форма дизайнерського мислення має бути інтегрована у технологічний та комунікативний компоненти ІКК. Здатність ефективно комунікувати з ШІ-системами через текстові промпти є специфічним дизайнерським умінням, що вимагає глибокого розуміння візуальної культури, стилістичних конвенцій та проєктної логіки.

\textit{По-третє}, проблеми упередженості ШІ, академічної доброчесності та авторського права обґрунтовують необхідність посилення рефлексивного компоненту ІКК. Критична рефлексія щодо етичних, правових та культурних аспектів використання ШІ є невід'ємною частиною сучасної дизайнерської професійності.

\textit{По-четверте}, теоретична недорозвиненість поля (73,7\,\% атеоретичних досліджень) та географічна лакуна (відсутність досліджень з українського контексту) обґрунтовують актуальність та оригінальність даного дисертаційного дослідження. Розробка теоретично обґрунтованої, емпірично перевіреної моделі формування ІКК для підготовки українських дизайнерів заповнює суттєву прогалину у наявній науковій літературі.

\textit{По-п'яте}, виявлений <<компетентнісний розрив>> --- розбіжність між швидкістю впровадження ШІ та розвитком педагогічних стратегій --- підтверджує необхідність системного, а не стихійного підходу до формування ІКК. Результати огляду свідчать, що ad hoc інтеграція ШІ без належного педагогічного обґрунтування не забезпечує ефективного розвитку компетентностей.

Таким чином, скопінг-огляд 156 досліджень засвідчив, що дослідницьке поле ШІ в дизайн-освіті перебуває у стані вибухового зростання (6,4-кратне збільшення після ChatGPT), але водночас характеризується суттєвими теоретичними (73,7\,\% атеоретичних), методологічними (53,8\,\% нижчої якості), дисциплінарними (графічний дизайн --- 1,9\,\%) та географічними лакунами. Ідентифіковані тематичні кластери --- креативність (35,9\,\%), оцінювання (27,6\,\%), дизайн навчальних програм (22,4\,\%) --- визначають ключові зони, у яких формування ІКК має забезпечити адекватну відповідь на виклики ШІ-трансформації. Парадигмальний зсув <<від виробника до куратора>> та виникнення промпт-інженерії як нової компетентності безпосередньо обґрунтовують структуру та зміст ІКК, що буде деталізовано у подальших підрозділах дисертації.
